\section{Labyrinth}
\label{sec:cses-labyrinth}

\problemstatement{Given a grid with a start \texttt{A}, a target
\texttt{B}, floor \texttt{.}, and walls \texttt{\#}, find a shortest path
from \texttt{A} to \texttt{B} moving in the four directions. Output
\texttt{YES}, its length, and the sequence of moves (\texttt{U/D/L/R}), or
\texttt{NO}.}

\begin{patternbox}
A grid where each move costs one step is an unweighted graph, so
\textbf{BFS} finds the shortest path, and remembering \emph{how} each cell
was first reached (which move and from where) reconstructs the actual path.
\end{patternbox}

\subsection*{BFS plus per-cell parent}

Run BFS from \texttt{A}. The first time BFS reaches a cell, it does so by a
shortest path, so record the move (\texttt{U/D/L/R}) used to enter it and
its predecessor. When \texttt{B} is reached, walk the parent pointers back
to \texttt{A}, collecting the moves, then reverse them. If BFS finishes
without reaching \texttt{B}, output \texttt{NO}.

\begin{ideabox}[Store the Move That First Reached Each Cell]
Distances alone say how long the path is; storing the direction taken into
each cell lets you replay the path. Because BFS's first arrival is optimal,
these stored moves spell out one shortest route.
\end{ideabox}

\subsection*{Algorithm}

\begin{enumerate}
  \item BFS from \texttt{A}; for each newly visited cell store the move
        character that entered it.
  \item On reaching \texttt{B}, backtrack via parents, collect moves,
        reverse.
  \item Output \texttt{YES}, the length, and the move string; else
        \texttt{NO}.
\end{enumerate}

\subsection*{Dry run}

\dryrun{\texttt{A} at $(0,0)$, \texttt{B} at $(0,2)$, top row open.}

\begin{itemize}
  \item BFS reaches $(0,1)$ via \texttt{R}, then $(0,2)$ via \texttt{R}.
  \item Backtrack from \texttt{B}: moves \texttt{R},\texttt{R}; reversed
        gives \texttt{RR}, length $2$.
\end{itemize}

\subsection*{C++ implementation}

\begin{lstlisting}
#include <bits/stdc++.h>
using namespace std;

int main() {
    int n, m;
    cin >> n >> m;
    vector<string> g(n);
    for (auto& row : g) cin >> row;

    int sr, sc, er, ec;
    for (int i = 0; i < n; i++)
        for (int j = 0; j < m; j++) {
            if (g[i][j] == 'A') { sr = i; sc = j; }
            if (g[i][j] == 'B') { er = i; ec = j; }
        }

    vector<vector<char>> move(n, vector<char>(m, 0));
    vector<vector<bool>> vis(n, vector<bool>(m, false));
    int dr[] = {-1,1,0,0}, dc[] = {0,0,-1,1};
    char mv[] = {'U','D','L','R'};

    queue<pair<int,int>> q;
    q.push({sr, sc}); vis[sr][sc] = true;
    while (!q.empty()) {
        auto [r, c] = q.front(); q.pop();
        for (int k = 0; k < 4; k++) {
            int nr = r + dr[k], nc = c + dc[k];
            if (nr < 0 || nr >= n || nc < 0 || nc >= m) continue;
            if (g[nr][nc] == '#' || vis[nr][nc]) continue;
            vis[nr][nc] = true;
            move[nr][nc] = mv[k];
            q.push({nr, nc});
        }
    }

    if (!vis[er][ec]) { cout << "NO\n"; return 0; }
    string path;
    int r = er, c = ec;
    while (!(r == sr && c == sc)) {
        char mvc = move[r][c];
        path += mvc;
        if (mvc == 'U') r++; else if (mvc == 'D') r--;
        else if (mvc == 'L') c++; else c--;      // step back
    }
    reverse(path.begin(), path.end());
    cout << "YES\n" << path.size() << "\n" << path << "\n";
    return 0;
}
\end{lstlisting}

\begin{complexitybox}
\textbf{Time Complexity.} $O(nm)$ -- one BFS over the grid.
\textbf{Space Complexity.} $O(nm)$ for visited and move grids.
\end{complexitybox}

\begin{takeawaybox}
Shortest path on an unweighted grid is BFS; storing the entering move per
cell turns the distance into a printable route. The backtrack-and-reverse
pattern reappears in every ``print the path'' problem, including Message
Route next.
\end{takeawaybox}
